% ============================================================
% CAPITOLO 4 - STATI E RIVOLTI DELLA TRIADE
% ============================================================
\chapter{Stati e rivolti della triade}
\label{chap:rivolti}

Le triadi studiate nel Capitolo \ref{chap:triadi} avevano la fondamentale al basso, cioè erano in \textbf{stato fondamentale}. Le tre note si possono però disporre in altro modo, ponendo al basso la terza o la quinta: sono i \textbf{rivolti}.

\section{Basso fondamentale e basso reale}
\label{subsec:basso_fondamentale_reale}

Prima dei rivolti serve una distinzione terminologica.

\begin{definizione}[Basso fondamentale e basso reale]
	\label{def:basso_fondamentale_reale}
	In un accordo si distinguono due concetti:
	\begin{itemize}
		\item \textbf{Basso fondamentale:} la nota che dà il nome all'accordo e su cui è costruito per terze sovrapposte. È un concetto teorico, indipendente dalla disposizione delle note.
		\item \textbf{Basso reale:} la nota più grave effettivamente suonata.
	\end{itemize}
	Se basso reale e basso fondamentale coincidono l'accordo è in stato fondamentale, altrimenti è rivoltato.
\end{definizione}

\begin{esempio}[Basso fondamentale e basso reale in Do maggiore]
	\label{ex:basso_fond_reale}
	Lo stesso accordo di Do maggiore, con basso fondamentale Do, su tre bassi reali diversi.
	\musica{es_cap02_03}
	L'accordo è sempre «Do maggiore»; cambia soltanto il basso reale.
\end{esempio}

\section{Lo stato di sesta}
\label{subsec:primo_rivolto}

Il primo rivolto si ottiene ponendo la \textbf{terza} dell'accordo al basso.

\begin{definizione}[Primo rivolto delle triadi]
	\label{def:primo_rivolto}
	Un accordo è in \textbf{primo rivolto} quando la sua terza si trova al basso. Si indica con la cifra \textbf{6} (o 6/3), perché sul basso la fondamentale sta a distanza di sesta e la quinta a distanza di terza.
\end{definizione}

\begin{esempio}[Una triade in primo rivolto]
	\label{ex:primo_rivolto_triadi}
	\musica{es_cap02_04}
\end{esempio}

\section{Lo stato di quarta e sesta}
\label{subsec:secondo_rivolto}

Il secondo rivolto si ottiene ponendo la \textbf{quinta} dell'accordo al basso. Qui se ne dà solo la costruzione; quando e come usarlo si vede nel Capitolo \ref{chap:quarta_sesta}.

\begin{definizione}[Secondo rivolto delle triadi]
	\label{def:secondo_rivolto}
	Un accordo è in \textbf{secondo rivolto} quando la sua quinta si trova al basso. Si indica con la cifra \textbf{6/4}, perché sul basso la fondamentale sta a distanza di quarta e la terza a distanza di sesta.
\end{definizione}

\begin{esempio}[Una triade in secondo rivolto]
	\label{ex:secondo_rivolto_triadi}
	\musica{es_cap02_05}
\end{esempio}

\section{La numerica delle triadi}
\label{subsec:cifratura_accordi}

Lo stato di un accordo si scrive sotto il suo grado con una \textbf{numerica} (o cifratura). Per capire da dove nasce basta guardare l'accordo dal basso.

\begin{definizione}[Numerica di un accordo]
	\label{def:numerica_accordo}
	La \textbf{numerica completa} di un accordo è l'elenco degli intervalli che le altre note formano con il \textbf{basso reale}, contati dal basso senza badare alle ottave e scritti in colonna, dal più grande in alto al più piccolo in basso. Non dipende dalla tonalità né dal grado: dipende solo da quale nota dell'accordo sta al basso.
\end{definizione}

Poiché alcuni numeri sono sottintesi, si usa una \textbf{numerica sintetica}, che conserva le sole cifre necessarie a distinguere gli stati. La tabella confronta le due forme per le triadi e per le settime (queste ultime si studiano dal Capitolo \ref{chap:settima_dominante}); gli esempi usano l'accordo di Do maggiore e la settima di dominante Sol-Si-Re-Fa.

\begin{center}
\small\renewcommand{\arraystretch}{1.25}
\begin{tabular}{llclc}
\toprule
\textbf{Stato} & \textbf{Basso} & \textbf{Completa} & \textbf{Sintesi possibili} & \textbf{Usata} \\
\midrule
Triade, fondamentale & Do (fondamentale) & 5 3 & 5 3, 5, \textit{nessuna} & --- \\
Triade, primo rivolto & Mi (terza) & 6 3 & 6 3, 6 & 6 \\
Triade, secondo rivolto & Sol (quinta) & 6 4 & 6 4 & \cifre{6}{4} \\
\midrule
Settima, fondamentale & Sol (fondamentale) & 7 5 3 & 7 5 3, 7 & 7 \\
Settima, primo rivolto & Si (terza) & 6 5 3 & 6 5 3, 6 5, 5 & \cifre{6}{5} \\
Settima, secondo rivolto & Re (quinta) & 6 4 3 & 6 4 3, 4 3 & \cifre{4}{3} \\
Settima, terzo rivolto & Fa (settima) & 6 4 2 & 6 4 2, 4 2, 2 & 2 \\
\bottomrule
\end{tabular}
\end{center}

La numerica completa si legge direttamente sulle note. Nel primo rivolto di Do maggiore (Mi-Sol-Do) sopra il Mi ci sono il Sol, a distanza di terza, e il Do, a distanza di sesta. Nel secondo rivolto (Sol-Do-Mi) sopra il Sol ci sono il Do (quarta) e il Mi (sesta). Nella settima di dominante con il Fa al basso, sopra il Fa ci sono il Sol (seconda), il Si (quarta) e il Re (sesta): numerica completa $6\,4\,2$, sintetica 2.

\begin{regola}[Numerica usata nel manuale]
	\label{reg:numerica_manuale}
	Nel manuale ogni accordo è indicato dal suo \textbf{grado} (numero romano), con sotto la \textbf{numerica sintetica}:
	\begin{itemize}
		\item \textbf{triadi:} stato fondamentale, nessuna cifra; primo rivolto 6; secondo rivolto 6 e 4 (la 6 sopra la 4);
		\item \textbf{settime:} stato fondamentale 7 (la cifra distingue la settima dalla triade); primo rivolto 6 e 5; secondo rivolto 4 e 3; terzo rivolto 2.
	\end{itemize}
	Nelle figure il grado è in maiuscolo per gli accordi maggiori (con $+$ se eccedenti) e in minuscolo per quelli minori e diminuiti (con $^\circ$ la triade e la settima diminuita, con ø la settima semidiminuita): per esempio I, ii, vii$^\circ$, viiø. Nel testo il grado si scrive di norma in maiuscolo. Il grado sta sopra e la numerica sotto, in colonna (per esempio il grado V con sotto 6 e 5). Nel testo corrente, per non alterare l'interlinea, la stessa sigla è scritta in forma compatta, con la prima cifra in alto e la seconda in basso: \sigla{I}{6}{}, \sigla{I}{6}{4}, \sigla{V}{7}{}, \sigla{V}{6}{5}, \sigla{V}{4}{3}, \sigla{V}{2}{}.
\end{regola}

La numerica sintetica va distinta dalla numerica del basso numerato, in cui si aggiungono altri numeri (ritardi come 4-3 o 9-8) e alterazioni ($\sharp$, $\flat$, $\natural$): se ne parla nel Capitolo \ref{chap:metodo_lavoro}.
